% part: intuitionistic-logic
% chapter: introduction
% section: syntax

\documentclass[../../../include/open-logic-section]{subfiles}

\begin{document}

\olfileid{int}{int}{syn}

\olsection{直觉主义逻辑的语法}

直觉主义逻辑的语法与命题逻辑的语法相同。
在经典命题逻辑中，可以用其他联结词来定义某些联结词，
例如，可以用 $\lnot !A \lor !B$ 来定义 $!A \lif !B$，
或者用 $\lnot(\lnot !A \land \lnot !B)$ 来定义 $!A \lor !B$。
因此，经典逻辑的表述常常将某些联结词作为上述定义的缩写引入。
直觉主义逻辑则不然，但有两个例外：
$\lnot !A$ 可以——而且经常是——被定义为 $!A \lif \lfalse$ 的缩写。
当然，这样一来 $\lfalse$ 本身就不能再被定义了！
此外，$!A \liff !B$ 可以像在经典逻辑中一样，定义为 $(!A \lif !B) \land (!B \lif
!A)$。

命题直觉主义逻辑的公式借助\emph{逻辑联结词}由\emph{命题变元}和命题常项~$\lfalse$
构建而成。具体包括：

\begin{enumerate}
\item 一个可数无穷的命题变元集~$\PVar$：$\Obj p_0$, $\Obj p_1$, \dots
\item 表示假的命题常项~$\lfalse$。
\item 逻辑联结词：$\land$（合取）、$\lor$（析取）、$\lif$（条件句）。
\item 标点符号：(, ) 以及逗号。
\end{enumerate}

\begin{defn}[公式]
\ollabel{defn:formulas}
命题直觉主义逻辑的\emph{公式}的集合~$\Frm[L_0]$ 归纳定义如下：
\begin{enumerate}
\item $\lfalse$ 是原子公式。
\item 每个命题变元~$\Obj p_i$ 都是原子公式。
\item 若 $!A$ 和 $!B$ 是公式，则 $(!A \land !B)$ 是公式。
\item 若 $!A$ 和 $!B$ 是公式，则 $(!A \lor !B)$ 是公式。
\item 若 $!A$ 和 $!B$ 是公式，则 $(!A \lif !B)$ 是公式。
\tagitem{limitClause}{除此之外均非公式。}{}
\end{enumerate}
\end{defn}

除了上面引入的原始联结词，我们还使用以下\emph{被定义的}符号：
$\lnot$（否定）和 $\liff$（双条件句）。
使用这些被定义的算子构造的公式应按以下方式理解：

\begin{enumerate}
\item $\lnot !A$ 是 $!A \lif \lfalse$ 的缩写。
\item $!A \liff !B$ 是 $(!A \lif !B) \land (!B
  \lif !A)$ 的缩写。
\end{enumerate}

尽管 $\lnot$ 在正式处理中被视为缩写，我们有时仍会像对待原始符号一样，为~$\lnot$ 给出显式的规则和定义条款。这主要是为了方便陈述练习题。

\end{document}